\documentclass[11pt]{article}
\usepackage[margin=1in]{geometry}
\usepackage{amsmath,amssymb,amsthm,hyperref,enumitem}
\hypersetup{colorlinks=true,linkcolor=blue,urlcolor=blue}
\newtheorem{theorem}{Theorem}
\newcommand{\tr}{\operatorname{tr}}
\title{The next-to-leading trace coefficient\\of a two-peg web world}
\author{Principia Math In Progress}
\date{22 September 2026}
\begin{document}
\maketitle
\noindent\textbf{MathDB 332639. Full all-size enumeration, with the source index corrected.}

\section*{Task 1: the intended trace enumeration}
Following Dukes and White [1, Section~7], let \(W_n\) be the web world
on two pegs with \(n\) edges. Its diagrams are
\[
 D_\pi=\{(1,2,i,\pi_i):1\leq i\leq n\},\qquad \pi\in S_n.
\]
A \(k\)-colouring is a surjection from the edges to \([k]\).
Reconstruction retains the relative endpoint orders within each
colour and stacks the colour classes in increasing colour order.
The coefficient \([x^k]M^{(W_n)}_{D_\pi,D_\sigma}(x)\) counts
colourings of \(D_\pi\) whose reconstruction is \(D_\sigma\).
Thus the trace counts pairs consisting of a diagram and a
colouring reconstructing that same diagram.

Conjecture~29 asks for the next-to-leading coefficient
\[
 a_n=[x^{n-1}]\tr M^{(W_n)}(x),\qquad n\geq2.
\]
The source's Figure~7 gives \(2,8,42,264,1920,15840\) for
\(n=2,\ldots,7\) and identifies them with the sequence A121635.
Its displayed formula has a one-place index shift. The theorem
below proves the intended enumeration for every \(n\), including
the precise relation to that sequence.

\begin{theorem}
For every integer \(n\geq2\),
\[
 \boxed{\displaystyle
 a_n=(n-1)!\left(\binom n2+1\right)
      =\frac{(n-1)!(n^2-n+2)}2.}                         \tag{1}
\]
More precisely, if \(\mathrm{id}_n=(1,2,\ldots,n)\) and
\(\tau_j=(1,\ldots,j-1,j+1,j,j+2,\ldots,n)\), then
\[
 [x^{n-1}]M^{(W_n)}_{D_\pi,D_\pi}(x)=
 \begin{cases}
  \binom n2(n-1)!, & \pi=\mathrm{id}_n,\\
  (n-2)!, & \pi=\tau_j,\quad 1\leq j<n,\\
  0, & \text{otherwise}.
 \end{cases}                                            \tag{2}
\]
\end{theorem}

\begin{proof}
An onto colouring of \(n\) edges with \(n-1\) colours has
exactly one class of size two; every other class is a singleton.
A singleton reconstructs to one edge. The two-edge class
reconstructs either to the permutation \(12\), when its endpoints
have the same relative order on the two pegs, or to \(21\),
when their relative orders are opposite.

Stacking these blocks gives either \(\mathrm{id}_n\) or a single
adjacent transposition. Hence no other permutation can be equal
to its reconstruction, proving the last case of (2).

For \(D_{\mathrm{id}_n}\), every pair has the same endpoint order,
so every surjective colouring reconstructs the identity.
There are \(\binom n2\) choices of the two-edge class, followed
by \((n-1)!\) bijections from the resulting \(n-1\) classes
to the labelled colours. This is the first case of (2).

For \(D_{\tau_j}\), the unique pair with opposite endpoint orders
consists of the edges whose left endpoints are \(j,j+1\).
They must be the two-edge class: any other choice reconstructs
the identity. If that class has colour \(h\), exactly \(h-1\)
singleton blocks precede it, so the reconstructed transposition
is \(\tau_h\). Thus reconstruction equals \(\tau_j\) exactly
when the two-edge class has colour \(j\).
The other \(n-2\) edges may receive the remaining \(n-2\) colours
in any bijective order, giving \((n-2)!\) colourings.
This proves the second case of (2), including \(n=2\) with \(0!=1\).

There are \(n-1\) adjacent transpositions. Summing (2) gives
\[
 a_n=\binom n2(n-1)!+(n-1)(n-2)!
    =(n-1)!\left(\binom n2+1\right),
\]
as claimed.
\end{proof}

\section*{The index and the sequence identification}
The known formula recorded for the polyomino sequence [2] is
\[
 B_m=\frac{(m-2)!(m^2-3m+4)}2\quad(m\geq2),\qquad B_1=1.
\]
Equation (1) proves \(a_n=B_{n+1}\) for every \(n\geq2\).
Equivalently, for every \(m\geq3\),
\[
 [x^{m-2}]\tr M^{(W_{m-1})}(x)
       =\frac{(m-2)!(m^2-3m+4)}2.                         \tag{3}
\]
The unshifted printed formula gives \(1\) at \(n=2\),
whereas the source itself gives \(2\). Merely observing that
discrepancy would not solve the intended counting question.
The proof above instead determines all diagonal contributions
and proves the source's intended sequence identification at
every size. The polyomino enumeration formula is attributed
to the existing sequence record; no new polyomino theorem is claimed.

This proof uses the definition of reconstruction directly.
Theorem~28 of the source gives a broader reduction to words;
the general enumeration in Problem~25 remains outside this result.

\section*{Tasks 2--5: proposed mathematical follow-ups}
\begin{enumerate}[label=\textbf{Task \arabic*:},start=2,leftmargin=*]
\item Determine a closed formula for \([x^{n-2}]\tr M^{(W_n)}(x)\)
for all \(n\geq3\), treating the cases of one three-edge class
and two two-edge classes.
\item For each fixed deficit \(d\), express
\([x^{n-d}]\tr M^{(W_n)}(x)\) by a finite classification of
non-singleton blocks, and determine its dependence on \(n\).
\item Construct an explicit bijection from the pairs counted
in (1) to deco polyominoes of height \(n+1\) with no two-cell
column starting at level zero, explaining the index shift.
\item Determine the next trace coefficients refined by the
number of inversions of the reconstructed permutation, and
derive recurrences for the resulting bivariate polynomials.
\end{enumerate}

\begin{thebibliography}{9}
\bibitem{dw}
M.~Dukes and C.~D.~White, \emph{Web matrices: structural properties
and generating combinatorial identities},
Electronic Journal of Combinatorics \textbf{23}(1) (2016), P1.45,
Section~7, pp.~21--23.
\href{https://emis.de/ft/29232}{Final journal version};
\href{https://arxiv.org/abs/1603.01589}{arXiv:1603.01589v1}.
\bibitem{oeis}
E.~Deutsch and subsequent contributors,
\href{https://oeis.org/A121635}{OEIS A121635},
sequence and formula for deco polyominoes; accessed
22 September 2026. The web-matrix references in this entry
are the conjectural papers in [1].
\end{thebibliography}
\end{document}
